\section{Diagnosis}
\label{sec:diagnosis}

\noindent
The symptoms gathered in \xref{sec:symptomatology} share a shape, and a shape that recurs across scales otherwise unrelated is evidence of a common cause. This section names that cause. It does so in the manner the foundational paper established for such work~\citep{wanhong_grb}: diagnosis here is not a preliminary to be hurried through on the way to a remedy, but a substantial undertaking in its own right, on the principle that the precision of a response is bounded by the precision of the diagnosis it answers. To see clearly how a gap has opened between the symbol and the thing it was to secure is already half of what is needed to respond to it; and so the section takes its time, and develops the diagnosis to a definite thesis, before any mechanism or remedy is proposed.

The foundational paper diagnosed a crisis of the symbol under modern conditions across several dimensions. The present paper does not recapitulate them all; it takes up the two that bear directly on trust\,---\,the crisis of authenticity and trust, and the overload of justice\,---\,and develops each beyond the point at which the foundational paper left it, weaving in three further of its dimensions where they bear on the mechanism, and adding one diagnostic dimension the foundational paper did not contain. It then sets the whole against the two older philosophical problems\,---\,Hume's and Kant's\,---\,that give the modern crisis its deeper form, and draws from the conjunction the thesis that organises the rest of the paper.

\subsection{The crisis of trust: from the failure of trust to the failure of its production}
\label{sub:diag-trust}

The first dimension the foundational paper diagnosed under this head is by now familiar from the symptomatology: a saturation of the human environment with symbols whose cost of production has fallen toward zero, whose provenance cannot be established, and which are as available to deception as to sincerity, with the rational consequence that the weight given to symbolic testimony as such declines~\citep{wanhong_grb}. The foundational paper established this much. What the present paper must add\,---\,for it is the precise point at which the diagnosis turns into the problem this paper exists to solve\,---\,is that the crisis has passed beyond the failure of particular instances of trust into a doubt about \emph{the producibility of trust at all}. It is one thing for a given promise to be disbelieved; it is another, and the deeper condition, for it to become unclear by what means \emph{any} promise could now be made believable. The symptomatology showed the apparatus of assurance failing at every scale; the diagnosis is that the failure is not of this or that instrument but of the question they all answer, so that the question itself\,---\,how is trust to be generated?\,---\,which in a healthier condition would not need to be asked, has become both urgent and obscure. This is why the mechanism of trust-generation, which a more confident age could take for granted and leave unexamined, must now be made the express object of inquiry.

\subsection{The overload of justice and the failure of force as a signal}
\label{sub:diag-justice}

The second dimension is the one the foundational paper named the overload of justice~\citep{wanhong_grb}. As relational and communal orders of trust have thinned, the burden of holding people to their dealings has fallen, by default, upon the juridical: where once a web of relationship and reputation and communal witness held a person to his word, now the law is asked to do it, and is asked to do it for a volume and a variety of dealings it was never sized to bear. The foundational paper diagnosed the resulting overload: the law extended into ever more of life, and at the same time hollowed from within, asked to underwrite an order it lacks the means to underwrite.

The present paper carries this diagnosis one decisive step. The overload is not merely quantitative, a matter of too many demands upon a finite institution; it is the visible form of a deeper failure, the failure of \emph{force as a signal}. The symptomatology recorded the sword that hangs but does not fall. The diagnosis of that symptom is this: the threat of compulsion produces compliance not by the mechanical certainty of the sanction\,---\,which, as the symptomatology showed, has become anything but certain\,---\,but by signalling to the parties that compliance is to be expected. Force, that is, operates as a substitute for expectation: it is one means among possible others of bringing a party to expect that an undertaking will be honoured. And a signal whose mechanical basis is known to have failed\,---\,a sword known not to fall\,---\,ceases to signal. The overload of justice is thus the macroscopic, institutional, observable face of the paper's central thesis, encountered here as a hard social fact before it is argued as a principle: that force was never the cause of the obligation, but only one device for the production of the expectation in which trust consists, and that the device is now failing on its own terms.

\subsection{Three further dimensions: cause, consequence, and precondition}
\label{sub:diag-three}

Three further of the foundational paper's dimensions enter the diagnosis, not as separate crises but as the cause, the consequence, and the precondition of the crisis of trust, and they are woven in here on that understanding rather than treated in their own right~\citep{wanhong_grb}.

The first is the \emph{generativity paradox}: the capacity of the present age's technologies to produce symbols without limit and at no cost. This is the \emph{cause}, at the level of material technique, of the saturation diagnosed above; it is why the forged signal has become not an occasional intrusion but the ambient condition, and why the production of a convincing token of sincerity can no longer, by itself, carry the weight it once did. The diagnosis of trust must be set upon this technical base, on pain of mistaking for a moral decline what is in the first instance a change in the cost of producing signs.

The second is the \emph{return of the Real}: the eruption, where the symbolic order fails to hold, of what the symbolic was meant to bind\,---\,in the form of a free-floating anxiety, a susceptibility to conspiracy, a drift toward the fundamentalist and the violent. This is the \emph{consequence}, at the level of the affective and the social, of a generalised failure of trust: when the symbolic testimony by which people orient themselves to one another can no longer be relied upon, what returns is not a neutral scepticism but a charged and disoriented seeking after some unfalsifiable ground. The diagnosis records this as the cost of leaving the crisis of trust unanswered.

The third is the \emph{collapse of intersubjectivity}, the degradation of the relation in which one meets another as a presence rather than as an object\,---\,what an earlier language opposed as the I--Thou to the I--It. This is the \emph{precondition} whose erosion most threatens the mechanism this paper will propose. For one of the channels by which trust is generated, as \xref{sec:mechanism} will show, requires that a real self-committing act be received \emph{by another who is present to receive it}; and where the relation has degraded to the merely instrumental, that channel is disabled at its root. The collapse of intersubjectivity is thus not one more symptom but a threat to the very possibility of the cure.

\subsection{The decoupling of credibility from enforceability}
\label{sub:diag-decoupling}

To the dimensions drawn from the foundational paper the present diagnosis adds one the foundational paper did not contain, and which is, the paper submits, the specifically modern hinge on which its central problem turns. It is the historical decoupling of the \emph{credibility} of a promise from its \emph{enforceability}.

In the pre-modern settings against which the modern condition is most sharply seen, the two were bound together and were scarcely distinguished. A promise was credible and enforceable at once, and by the same fact: the god who witnessed the oath was also the power that would avenge its breach; the lineage or the village before whom the word was given was also the body that would exact the cost of breaking it; to be seen to promise was, in the same act, to be bound to perform, because the seeing and the binding were done by one and the same watching community. Credibility and enforceability were two aspects of a single social fact, and no one had occasion to prise them apart.

Modernity prised them apart. It assigned enforceability to a specialised institution, the law, which would compel performance through an abstract and impersonal machinery; and it left credibility\,---\,the quality by which a promise is believed in the first place, before and apart from any question of compulsion\,---\,without a corresponding institution to house it. So long as the law's machinery of enforcement was robust, this division of labour was invisible, because enforceability could be relied upon to stand in for credibility: one did not need to find a promise credible if one could be assured it was enforceable. But the symptomatology has shown the enforcing machinery failing; and as it fails, the division of labour is exposed, and credibility is revealed for what modernity had quietly made it\,---\,an orphan, a quality with no institution to underwrite it, hanging unsupported now that the institution which had stood in for it can no longer bear the weight.

This is why the paper's central problem is urgent \emph{now}, and was not, in this form, urgent before. The question of how a promise becomes credible in the absence of enforcement is not a perennial philosophical puzzle that happens to have been posed again; it is the question a specific historical arrangement has forced into the open, by first making credibility depend upon enforceability and then allowing enforceability to fail. The orphaning of credibility is the diagnosis that pins the paper to its moment, and it is offered as the diagnostic contribution most particularly the paper's own.

\subsection{The two older problems: Hume and Kant}
\label{sub:diag-humekant}

Beneath the modern crisis lie two older problems, which give it its deeper philosophical form and which the diagnosis must reach if the thesis it arrives at is to be more than a sociological observation.

The first is Hume's. The expectation that a person who has kept his word will keep it again, like every expectation that the future will resemble the past, rests upon induction, and induction, as Hume showed, has no foundation in reason: no accumulation of past performance \emph{entails} future performance, and the inference from the one to the other, however indispensable, cannot be justified as a matter of logic~\citep{hume2000}. This is the problem that the evidentialist's reliance upon track record can never wholly escape, and it sets a limit on what any quantity of evidence can establish: at the end of the longest record of kept promises there remains a gap, across which no evidence can carry one, between what has been and what will be. The diagnosis records this as the deep structural reason why trust can never be reduced to the accumulation of evidence: the gap Hume opened is constitutive, and every account of trust must say what it does about that gap rather than pretend to have closed it.

The second is Kant's. Where Hume's problem is that the inference to future conduct cannot be grounded in past fact, Kant's bears on the standing of the one who must nonetheless be relied upon. To treat a person solely as a system whose future behaviour is to be predicted from evidence and secured by incentive is to treat him, in the terms of the second formulation of the categorical imperative, merely as a means, as an object to be managed rather than an end in himself~\citep{kant1998}. But a promise, in its full sense, is addressed to a free agent who binds himself, and is received from one; the one who promises is not predicting his own future conduct as one predicts the weather, but committing himself to it as only a free being can. The diagnosis records this as the reason the problem of trust cannot be fully handed over to prediction and incentive without changing its object: to model the promisor merely as a mechanism to be managed is to have stopped speaking of promising at all. The wager that the other will keep faith is, in part, a wager on his freedom\,---\,and this, as \xref{sec:freedom} will develop, is not a softness in the account but the very opening in which everything the paper most wants to say becomes possible.

\subsection{The thesis: force is not the cause of trust}
\label{sub:diag-thesis}

The dimensions assembled converge on a single thesis, which the rest of the paper exists to develop, to ground, and to apply. It can be stated plainly. \emb{Even where the coercive guarantee of the Other is fully present, trust under modern conditions has broadly declined; therefore force is not the cause of trust.} If the sword were the cause of the expectation, then where the sword stands ready the expectation should stand firm; but the symptomatology shows trust eroding precisely where the apparatus of compulsion is most fully elaborated, and the diagnosis of the overload of justice shows why: the sword secures expectation only by signalling, and a signal whose basis is known to have failed no longer signals, however sharp the blade remains. Force, then, is at most a contingent and now-failing device for the production of an expectation that is the real thing; it is not that expectation's cause.

The thesis turns on a distinction the rest of the paper will hold to throughout: between two senses of \emph{cause}. There is the mechanical cause, which operates as one billiard ball operates upon another, by the transmission of a compulsion that leaves the affected body no part to play; and there is the inferential cause, which operates by giving a mind a ground on which to form an expectation, an expectation the mind itself produces and could, being free, have withheld. Force, where it works, works mechanically, or aspires to: it would move the promisor as a body is moved. But trust is not a body's motion; it is a mind's expectation, formed inferentially, on grounds. To say that force is not the cause of trust is to say that the mechanical cause is not the cause of the inferential effect\,---\,that the production of an expectation in a free mind belongs to a different order of causation than the compulsion of a body, and cannot be accomplished by it. The whole positive account of the paper follows from taking this distinction seriously: if trust is an inferential effect, then the question of how it is generated is the question of how the right grounds are supplied to the mind that forms it\,---\,and force, it turns out, is neither the only such ground nor, any longer, a working one.

With this the diagnosis is complete. The symptom was a gap, recurring across every scale, between the apparatus of assurance and the trust it failed to secure. The diagnosis is that the apparatus failed because it was built on a mistake about its own working\,---\,the mistake of taking force, the mechanical cause, for the cause of an effect that is inferential and free\,---\,and that the crisis of the present is the coming-due of that mistake, as the one device the mistake had relied upon, enforcement, fails on its own terms and leaves the credibility it had stood in for without support. What remains is to show, positively, how the inferential effect is in fact produced: by what mechanism a free mind, given the right grounds, comes to expect that a promise will be kept. To that the paper now turns, beginning where the failed device itself began, in the philosophy of law.
